\clearpage
\fancyhead[L]{\small THEORIE DER INTERPOLATION}
\pdfbookmark[0]{Theorie der Interpolation nach einer neuen Methode behandelt}{interpolation-nachlass}
\begingroup
\renewcommand{\arraystretch}{1.35}
\SourceLeaf{274}
\begin{center}\textbf{NACHLASS.}\par\bigskip
\textbf{THEORIE DER INTERPOLATION}\par\medskip
\textbf{NACH EINER NEUEN METHODE BEHANDELT.}\par\medskip
\rule{.15\linewidth}{.4pt}\end{center}
\InterpolationArticle{1}
\textsc{Aufgabe.}\quad\textit{Die Summe der Reihe zu finden}
\[
\begin{aligned}
&\frac{a^n}{(a-b)(a-c)(a-d)(a-e)\ldots}\\[5pt]
&+\frac{b^n}{(b-a)(b-c)(b-d)(b-e)\ldots}\\
&+\frac{c^n}{(c-a)(c-b)(c-d)(c-e)\ldots}\\[5pt]
&+\frac{d^n}{(d-a)(d-b)(d-c)(d-e)\ldots}\\
&+\frac{e^n}{(e-a)(e-b)(e-c)(e-d)\ldots}+\text{usw.}
\end{aligned}
\]
\textit{wobei $a,b,c,d,e$ $m$ verschiedene Größen sind und $n$ eine beliebige ganze Zahl ist, positiv, negativ oder auch 0.}

\textit{Lösung.}\quad Der Kürze halber setze man
\[
\begin{aligned}
\frac1{(a-b)(a-c)(a-d)(a-e)\ldots}&=\alpha\\
\frac1{(b-a)(b-c)(b-d)(b-e)\ldots}&=\beta\\
\frac1{(c-a)(c-b)(c-d)(c-e)\ldots}&=\gamma\\
\frac1{(d-a)(d-b)(d-c)(d-e)\ldots}&=\delta,\ \text{usw.}
\end{aligned}
\]
sodass die gesuchte Summe, die wir mit $S^n$ bezeichnen wollen, $=\alpha a^n+\beta b^n+\gamma c^n+\delta d^n+$ usw. wird. Bezeichnet $x$ eine unbestimmte Größe, so ergibt sich offenbar aus der Entwicklung der Summe
\SourceLeaf{275}
\[
P=\frac\alpha{1-ax}+\frac\beta{1-bx}+\frac\gamma{1-cx}+\frac\delta{1-dx}+\text{usw.}
\]
in eine Reihe nach aufsteigenden Potenzen von $x$
\[
S^0+S^1x+S^2xx+S^3x^3+\text{usw.}\quad\text{ins Unendliche}
\]
Man setze $(1-ax)(1-bx)(1-cx)(1-dx)\ldots=Q$; dann ist $Q$ eine ganzrationale Funktion der Unbestimmten $x$, die bis zum $m^{\mathrm{ten}}$ Grad reicht; $PQ$ wird dagegen eine ganzrationale Funktion vom $m-1^{\mathrm{ten}}$ Grad, nämlich $=$
\[
\begin{aligned}
&\alpha(1-bx)(1-cx)(1-dx)\ldots\\
&+\beta(1-ax)(1-cx)(1-dx)\ldots\\
&+\gamma(1-ax)(1-bx)(1-dx)\ldots\\
&+\delta(1-ax)(1-bx)(1-cx)\ldots\\
&+\text{usw.}
\end{aligned}
\]
Bei näherer Betrachtung erkennt man, dass durch die Einsetzung $x=\frac1a$ alle Glieder außer dem ersten verschwinden, dieses aber übergeht in
\[
\alpha(1-\tfrac ba)(1-\tfrac ca)(1-\tfrac da)\ldots=\frac1{a^{m-1}}
\]
In gleicher Weise verschwinden durch die Einsetzung $x=\frac1b$ alle Glieder außer dem zweiten, das $=\frac1{b^{m-1}}$ wird. Ebenso geht durch die Einsetzungen $x=\frac1c$, $x=\frac1d$ usw. $PQ$ in $\frac1{c^{m-1}}$, $\frac1{d^{m-1}}$ usw. über. Daraus folgt, dass $PQ-x^{m-1}$ bei allen diesen Einsetzungen den Wert 0 erhält, was nur möglich ist, wenn es identisch $=0$ ist, also $PQ=x^{m-1}$; andernfalls hätte nämlich die Gleichung $PQ-x^{m-1}=0$, deren Grad den $m-1^{\mathrm{ten}}$ nicht übersteigt, $m$ verschiedene Wurzeln $\frac1a,\frac1b,\frac1c,\frac1d\ldots$.

Es sei nun
\[
Q=1-Ax+Bxx-Cx^3+Dx^4-\text{usw.}
\]
wobei $A$ die Summe der Größen $a,b,c,d\ldots$, $B$ die Summe ihrer Produkte aus je zweien, $C$ die Summe ihrer Produkte aus je dreien usw. bezeichnet. Da aus der Entwicklung des Bruchs $\frac{x^{m-1}}Q=P$
\[
S^0+S^1x+S^2xx+S^3x^3+\text{usw.}
\]
hervorgehen muss, gilt offenbar erstens $S^0=0$, $S^1=0$, $S^2=0$ usw. bis $S^{m-2}=0$\HistoricalNote{22}; dann
\SourceLeaf{276}
ist dagegen $S^{m-1}=1$, $S^m=A$, und schließlich werden die weiteren Glieder als Glieder einer rekurrenten Reihe nach folgendem Gesetz bestimmt:
\[
\begin{aligned}
S^m&=A\\
S^{m+1}&=AS^m-BS^{m-1}=AA-B\\
S^{m+2}&=AS^{m+1}-BS^m+CS^{m-1}=A^3-2BA+C\\
&\text{usw.}
\end{aligned}
\]
Hieraus lässt sich zwar leicht schließen, dass $S^{m+1}$ $=aa+bb+cc+\ldots+B$ ist, also die Summe der Quadrate zusammen mit der Summe aller Produkte aus je zwei verschiedenen der Größen $a,b,c,d\ldots$. Um aber deutlicher zu erkennen, wie die folgenden Glieder aus den Elementen $a,b,c,d\ldots$ gebildet werden, bemerken wir, dass $\frac1Q$ das Produkt der Reihen
\[
\begin{aligned}
&1+ax+aaxx+a^3x^3+\text{usw.}\\
&1+bx+bbxx+b^3x^3+\text{usw.}\\
&1+cx+ccxx+c^3x^3+\text{usw.}\\
&1+dx+ddxx+d^3x^3+\text{usw.}\\
&\text{usw.}
\end{aligned}
\]
ist. Dieses Produkt ist $=\Sigma a^\lambda b^\mu c^\nu\ldots\times x^{\lambda+\mu+\nu.\,.}$, wobei den Exponenten $\lambda,\mu,\nu\ldots$ alle ganzzahligen Werte von 0 bis ins Unendliche zu geben und diese auf alle möglichen Arten zu kombinieren sind. Um also in der Reihe, in die $\frac1Q$ entwickelt wird, den Koeffizienten des Gliedes mit $x^{n+1-m}$ zu erhalten, muss die Zahl $n+1-m$ auf alle möglichen Arten in $n$ ganzzahlige Teile $\lambda+\mu+\nu+.\,.$ zerlegt werden (wobei auch ein Teil 0 zugelassen ist), und zwar unter Zulassung aller Vertauschungen dieser Teile; dann sind die einzelnen Produkte $a^\lambda b^\mu c^\nu\ldots$ zu addieren. Ihre Summe ist der gesuchte Koeffizient und zugleich $=S^n$. Bei einiger Aufmerksamkeit erkennt man, dass die folgende Regel dieser vollständig gleichwertig ist: Aus den $m$ Größen $a,b,c,d\ldots$ sind alle Kombinationen von $n+1-m$ Elementen mit Wiederholungen zu bilden und jeweils als Produkte aufzufassen; deren Summe ist $=S^n$. Somit ist wie oben $S^{m+1}$ die Summe aller Produkte aus je zwei der Größen $a,b,c,d\ldots$, verschiedenen wie gleichen; $S^{m+2}$ die Summe aller Produkte aus je drei verschiedenen oder gleichen Größen usw.

Es bleibt uns nur noch zu zeigen, wie die Summe unserer Folge für negative Werte von $n$ zu bestimmen ist. Dazu bringen wir das erste Glied der Summe $S^{-n}$, nämlich $\frac{a^{-n}}{(a-b)(a-c)(a-d)\ldots}$, in folgende Form
\SourceLeaf{277}
\[
\pm\frac1{abcd\ldots}\times\frac{(\frac1a)^{m+n-2}}{(\frac1a-\frac1b)(\frac1a-\frac1c)(\frac1a-\frac1d)\ldots}
\]
\Needspace{15\baselineskip}
wobei das obere oder untere Zeichen zu wählen ist, je nachdem, ob $m$ ungerade oder gerade ist; dieselbe Umformung lässt sich auch auf die übrigen Glieder anwenden. Bezeichnet daher das Zeichen $T$ den Ausdruck, der aus $\frac1a,\frac1b,\frac1c,\frac1d.\,.$ ebenso entsteht wie $S$ aus $a,b,c,d\ldots$, so ist offenbar $S^{-n}=\pm\frac{T^{m+n-2}}{abcd\ldots}$. Auf diese Weise ist also dieser Fall auf den vorhergehenden zurückgeführt, und es ergibt sich
\[
S^{-1}=\pm\frac1{abcd\ldots},\qquad S^{-2}=\pm\frac{\frac1a+\frac1b+\frac1c+\frac1d+.\,.}{abcd\ldots},
\]
$S^{-3}$ ist gleich dem Produkt aus $\pm\frac1{abcd\ldots}$ und der Summe aller Produkte aus je zwei der Größen $\frac1a,\frac1b,\frac1c,\frac1d\ldots$, verschiedenen oder gleichen, usw.
\InterpolationArticle{2}
Wir wenden die vorhergehende Untersuchung auf den Fall an, dass den Größen $a,b,c,d$ imaginäre Werte gegeben werden. So gelangen wir sehr leicht zu einigen bemerkenswerten Beziehungen, deren Behandlung nach einer anderen Methode größere Schwierigkeiten bereiten würde. Es sei $E$ die Basis der natürlichen Logarithmen, $i$ die imaginäre Größe $\BareRoot{-1}$; anstelle der reellen Größen $a,b,c\ldots$ betrachten wir die imaginären Größen $E^{ia},E^{ib},E^{ic},E^{id}\ldots$ und $E^{-ia},E^{-ib},E^{-ic},E^{-id}\ldots$ und setzen
\[
\begin{aligned}
&\frac{E^{ina}}{(E^{ia}-E^{ib})(E^{ia}-E^{ic})(E^{ia}-E^{id})\ldots}\\
&+\frac{E^{inb}}{(E^{ib}-E^{ia})(E^{ib}-E^{ic})(E^{ib}-E^{id})\ldots}\\
&+\frac{E^{inc}}{(E^{ic}-E^{ia})(E^{ic}-E^{ib})(E^{ic}-E^{id})\ldots}\\
&+\frac{E^{ind}}{(E^{id}-E^{ia})(E^{id}-E^{ib})(E^{id}-E^{ic})\ldots}+\text{usw.}=S^n,\quad\text{und}
\end{aligned}
\]
\[
\begin{aligned}
&\frac{E^{-ina}}{(E^{-ia}-E^{-ib})(E^{-ia}-E^{-ic})(E^{-ia}-E^{-id})\ldots}\\
&+\frac{E^{-inb}}{(E^{-ib}-E^{-ia})(E^{-ib}-E^{-ic})(E^{-ib}-E^{-id})\ldots}\\
&+\frac{E^{-inc}}{(E^{-ic}-E^{-ia})(E^{-ic}-E^{-ib})(E^{-ic}-E^{-id})\ldots}\\
&+\frac{E^{-ind}}{(E^{-id}-E^{-ia})(E^{-id}-E^{-ib})(E^{-id}-E^{-ic})\ldots}+\text{usw.}=T^n
\end{aligned}
\]
\SourceLeaf{278}
\Needspace{14\baselineskip}
Bezeichnet man also die Anzahl der Größen $a,b,c,d\ldots$ mit $m$, so sind $S^0,S^1,S^2\ldots S^{m-2}$ sowie $T^0,T^1,T^2\ldots T^{m-2}$ sämtlich $=0$; ferner ist $S^{m-1}=T^{m-1}=1$; $S^m$ ist die Summe der Größen $E^{ia},E^{ib},E^{ic},E^{id}\ldots$; $S^{m+1}$ die Summe aller Produkte aus je zwei verschiedenen oder gleichen; $S^{m+2}$ die Summe der Produkte aus je dreien usw. Ebenso sind die Werte der Summen $T^m,T^{m+1},T^{m+2}$ usw. aus den Größen $E^{-ia},E^{-ib},E^{-ic},E^{-id}\ldots$ zu bilden.

\Needspace{21\baselineskip}
Da nun bekanntlich $E^{ix}+E^{-ix}=2\cos x$, $E^{ix}-E^{-ix}=2i\sin x$ gilt, erkennt man leicht, dass der Wert des Ausdrucks $\frac12(S^n+T^n)$, der für $n=0,1,2\ldots m-2$ $=0$, für $n=m-1$ dagegen $=1$ ist, für $n=m$
\[
=\cos a+\cos b+\cos c+\cos d\,.\,.\,.\,.
\]
wird; ebenso wird $\frac12(S^{m+1}+T^{m+1})$ die Summe der Kosinus aller Winkel, die durch Addition von \textit{je zwei} verschiedenen oder gleichen der Winkel $a,b,c,d.\,.$ entstehen; $\frac12(S^{m+2}+T^{m+2})$ wird die Summe der Kosinus aller Winkel, die durch Addition von \textit{je drei} derselben Winkel entstehen, usw. Ebenso ist $\frac{S^n-T^n}{2i}=0$ für $n=0,1,2\ldots m-1$; ferner $=\sin a+\sin b+\sin c+\sin d+.\,.$ für $n=m$; und ebenso ist $\frac{S^{m+1}-T^{m+1}}{2i}$ die Summe der Sinus aller Winkel, die durch Addition von je zwei verschiedenen oder gleichen der Winkel $a,b,c,d\ldots$ entstehen; $\frac{S^{m+2}-T^{m+2}}{2i}$ die Summe der Sinus aller Winkel, die durch Kombination von je drei derselben Winkel entstehen, usw.

Wir betrachten nun die Glieder der Summen $S^n,T^n$ näher. Es ist
\[
\begin{aligned}E^{ia}-E^{ib}&=E^{\frac12i(a+b)}(E^{\frac12i(a-b)}-E^{\frac12i(b-a)})\\
&=2iE^{\frac12i(a+b)}\sin\tfrac12(a-b)\end{aligned}
\]
ebenso
\[
E^{ia}-E^{ic}=2iE^{\frac12i(a+c)}\sin\tfrac12(a-c)\quad\text{usw.}
\]
Daher wird der Nenner des ersten Gliedes in $S^n$
\[
=(2i)^{m-1}E^{\frac12i(b+c+d\ldots+(m-1)a)}\sin\tfrac12(a-b)\sin\tfrac12(a-c)\sin\tfrac12(a-d)\ldots
\]
und, mit der Setzung $a+b+c+d+\ldots=s$, dieses Glied selbst
\[
=\frac{E^{i((n+1-\frac12m)a-\frac12s)}}{(2i)^{m-1}\sin\tfrac12(a-b)\sin\tfrac12(a-c)\sin\tfrac12(a-d)\ldots}
\]
In gleicher Weise gilt
\[
E^{-ia}-E^{-ib}=-2iE^{-\frac12i(a+b)}\sin\tfrac12(a-b)
\]
woraus schließlich das erste Glied der Summe $T^n$ hervorgeht als
\SourceLeaf{279}
\[
=\frac{\pm E^{-i((n+1-\frac12m)a-\frac12s)}}{(2i)^{m-1}\sin\tfrac12(a-b)\sin\tfrac12(a-c)\sin\tfrac12(a-d)\ldots}
\]
\Needspace{24\baselineskip}
wobei das obere oder untere Zeichen gilt, je nachdem, ob $m$ ungerade oder gerade ist. Indem wir dies zum ersten Glied der Summe $S^n$ addieren oder davon subtrahieren, schließen wir \textit{erstens}, dass bei ungeradem $m$ das erste Glied der Summe $\frac12(S^n+T^n)$
\[
=\frac{\cos((n+1-\frac12m)a-\frac12s)}{(2i)^{m-1}\sin\tfrac12(a-b)\sin\tfrac12(a-c)\sin\tfrac12(a-d)\ldots}
\]
und das erste Glied der Summe $\frac{S^n-T^n}{2i}$
\[
=\frac{\sin((n+1-\frac12m)a-\frac12s)}{(2i)^{m-1}\sin\tfrac12(a-b)\sin\tfrac12(a-c)\sin\tfrac12(a-d)\ldots}
\]
\textit{zweitens}, dass bei geradem $m$ das erste Glied der Summe $\frac12(S^n+T^n)$
\[
=\frac{\sin((n+1-\frac12m)a-\frac12s)}{2^{m-1}i^{m-2}\sin\tfrac12(a-b)\sin\tfrac12(a-c)\sin\tfrac12(a-d)\ldots}
\]
und das erste Glied der Summe $\frac{S^n-T^n}{2i}$
\[
=\frac{-\cos((n+1-\frac12m)a-\frac12s)}{2^{m-1}i^{m-2}\sin\tfrac12(a-b)\sin\tfrac12(a-c)\sin\tfrac12(a-d)\ldots}
\]
wird. Die folgenden Glieder der Ausdrücke $\frac12(S^n+T^n)$, $\frac{S^n-T^n}{2i}$ sind offenbar dem ersten vollständig analog und gehen daraus allein durch Vertauschung des Zeichens $a$ mit $b,c,d\ldots$ hervor. Bezeichnen wir nun mit $k$ einen beliebigen Winkel und setzen wir $n+1-\frac12m=\lambda$, so gelangen wir mithilfe der Formeln
\[
\begin{aligned}
\cos(\lambda a+k)&=\cos(k+\tfrac12s)\cos(\lambda a-\tfrac12s)\\&\qquad-\sin(k+\tfrac12s)\sin(\lambda a-\tfrac12s)\\
\sin(\lambda a+k)&=\cos(k+\tfrac12s)\sin(\lambda a-\tfrac12s)\\&\qquad+\sin(k+\tfrac12s)\cos(\lambda a-\tfrac12s)
\end{aligned}
\]
\Needspace{24\baselineskip}
ohne Schwierigkeit zur Summation folgender Reihen
\[
\begin{aligned}
&\frac{\cos(\lambda a+k)}{\sin\tfrac12(a-b)\sin\tfrac12(a-c)\sin\tfrac12(a-d)\ldots}\\
&+\frac{\cos(\lambda b+k)}{\sin\tfrac12(b-a)\sin\tfrac12(b-c)\sin\tfrac12(b-d)\ldots}\\
&+\frac{\cos(\lambda c+k)}{\sin\tfrac12(c-a)\sin\tfrac12(c-b)\sin\tfrac12(c-d)\ldots}\\
&+\frac{\cos(\lambda d+k)}{\sin\tfrac12(d-a)\sin\tfrac12(d-b)\sin\tfrac12(d-c)\ldots}+\text{usw.}=U^\lambda
\end{aligned}
\]
und
\SourceLeaf{280}
\[
\begin{aligned}
&\frac{\sin(\lambda a+k)}{\sin\tfrac12(a-b)\sin\tfrac12(a-c)\sin\tfrac12(a-d)\ldots}\\
&+\frac{\sin(\lambda b+k)}{\sin\tfrac12(b-a)\sin\tfrac12(b-c)\sin\tfrac12(b-d)\ldots}\\
&+\frac{\sin(\lambda c+k)}{\sin\tfrac12(c-a)\sin\tfrac12(c-b)\sin\tfrac12(c-d)\ldots}\\
&+\frac{\sin(\lambda d+k)}{\sin\tfrac12(d-a)\sin\tfrac12(d-b)\sin\tfrac12(d-c)\ldots}+\text{usw.}=V^\lambda.
\end{aligned}
\]
\Needspace{27\baselineskip}
\begin{center}\textit{Erster Fall: $m$ ist ungerade.}\end{center}
In diesem Fall ist für $\lambda$ keine ganze Zahl, sondern eine um den Bruch $\frac12$ vermehrte ganze Zahl erforderlich, also die Hälfte einer ungeraden Zahl. Die \textit{negativen} Werte von $\lambda$ brauchen nicht betrachtet zu werden, da $\cos(-\lambda a+k)=\cos(\lambda a-k)$ und $\sin(-\lambda a+k)=-\sin(\lambda a-k)$ gilt und daher die Summation für einen negativen Wert aus derjenigen für den entgegengesetzten positiven Wert allein durch Änderung von $k$ in $-k$ unmittelbar folgt. Diese Bemerkung gilt offenbar auch im folgenden Fall. Nun erkennt man leicht, dass für $n=0,1,\ldots m-2$, also für $\lambda=-\frac12m+1,-\frac12m+2,\ldots\frac12m-1$ oder, wenn wir die negativen Werte weglassen, für $\lambda=\frac12,\frac32,\frac52\ldots\frac12m-1$ sowohl $U^\lambda=0$ als auch $V^\lambda=0$ gilt; ferner gelten die Beziehungen
\[
\begin{aligned}
U^{\frac12m}&=(2i)^{m-1}\cos(\tfrac12s+k),\\V^{\frac12m}&=(2i)^{m-1}\sin(\tfrac12s+k)\\
U^{\frac12m+1}&=(2i)^{m-1}\{\cos(\tfrac12s+k+a)+\cos(\tfrac12s+k+b)\\&\qquad+\cos(\tfrac12s+k+c)\\
&\qquad+\cos(\tfrac12s+k+d)+\text{usw.}\}\\
V^{\frac12m+1}&=(2i)^{m-1}\{\sin(\tfrac12s+k+a)+\sin(\tfrac12s+k+b)\\&\qquad+\sin(\tfrac12s+k+c)\\
&\qquad+\sin(\tfrac12s+k+d)+\text{usw.}\}
\end{aligned}
\]
$U^{\frac12m+2}$ oder $V^{\frac12m+2}$ ist gleich dem Produkt aus $(2i)^{m-1}$ und der Summe der Kosinus oder Sinus aller Winkel, die entstehen, wenn $\frac12s+k$ zu je zwei verschiedenen oder gleichen der Winkel $a,b,c,d\ldots$ addiert wird; $U^{\frac12m+3}$ oder $V^{\frac12m+3}$ ist gleich dem Produkt aus $(2i)^{m-1}$ und der Summe der Kosinus oder Sinus aller Winkel, die entstehen, wenn $\frac12s+k$ zu je drei derselben Winkel addiert wird, usw. Übrigens braucht kaum daran erinnert zu werden, dass eine Potenz der imaginären Größe $i$ bei einem Exponenten, der das Doppelte einer geraden Zahl ist, $=+1$ wird, bei einem Exponenten, der das Doppelte einer ungeraden Zahl ist, dagegen $=-1$.
\Needspace{24\baselineskip}
\begin{center}\textit{Zweiter Fall: $m$ ist gerade.}\end{center}
In diesem Fall muss $\lambda$ eine ganze Zahl sein, und für $\lambda=0,1,2\ldots\frac12m-1$ gilt sowohl $U^\lambda=0$ als auch $V^\lambda=0$. Ferner gelten die Beziehungen
\SourceLeaf{281}
\[
\begin{aligned}
U^{\frac12m}&=2^{m-1}i^{m-2}\sin(\tfrac12s+k),\\
V^{\frac12m}&=-2^{m-1}i^{m-2}\cos(\tfrac12s+k);\\
U^{\frac12m+1}&=2^{m-1}i^{m-2}\{\sin(\tfrac12s+k+a)+\sin(\tfrac12s+k+b)\\&\qquad+\sin(\tfrac12s+k+c)\\
&\qquad+\sin(\tfrac12s+k+d)+\text{usw.}\}\\
V^{\frac12m+1}&=-2^{m-1}i^{m-2}\{\cos(\tfrac12s+k+a)+\cos(\tfrac12s+k+b)\\&\qquad+\cos(\tfrac12s+k+c)\\
&\qquad+\cos(\tfrac12s+k+d)+\text{usw.}\}
\end{aligned}
\]
$U^{\frac12m+2},U^{\frac12m+3}$ usw. sind gleich dem Produkt aus $2^{m-1}i^{m-2}$ und der Summe der Sinus aller Winkel, die entstehen, wenn $\frac12s+k$ zu je zwei, drei usw. verschiedenen oder gleichen der Winkel $a,b,c\ldots$ addiert wird; $V^{\frac12m+2},U^{\frac12m+3}$ usw. sind gleich dem Produkt aus $-2^{m-1}i^{m-2}$ und der Summe der Kosinus aller Winkel, die entstehen, wenn $\frac12s+k$ zu je zwei, drei usw. derselben Winkel addiert wird.
\endgroup
